\documentclass[10pt,a4paper]{amsart}
\usepackage[margin=19mm]{geometry}
\usepackage{amsmath,amssymb,mathtools}
\usepackage[scr=rsfs,cal=euler]{mathalfa}
\usepackage{xfrac}
\usepackage[unicode,hidelinks,bookmarks=false]{hyperref}
\newcommand{\ie}{i.e.\ }
\newcommand{\cf}{cf.\ }
\newcommand{\resp}{respectively\ }
\newcommand{\Subsection}{\subsection}
\providecommand{\nobreakdash}{\nobreak\hbox{-}}
\setlength{\emergencystretch}{2em}
\multlinegap0pt
\usepackage{bm}
\usepackage{bbm}
\usepackage{dsfont}
\usepackage{xspace}
\usepackage{cancel}
\usepackage{mathtools}
\usepackage{amsthm}
\makeatletter
\@ifundefined{theorem}{\newtheorem{theorem}{Theorem}}{}
\@ifundefined{proposition}{\newtheorem{proposition}[theorem]{Proposition}}{}
\makeatother
\makeatletter
\newcommand{\CC}{\mathcal{C}}
\newcommand{\N}{\mathbb{N}}
\newcommand{\PP}{\mathbb{P}}
\newcommand{\R}{\mathbb{R}}
\newcommand{\Set}{\mathfrak{S}}
\newcommand{\VV}{\mathcal{V}}
\newcommand{\XX}{\mathcal{X}}
\newcommand{\dd}{\mathrm{d}}
\expandafter\ifx\csname proof\endcsname\relax
\newenvironment{proof}{{\bf Proof.}}{\hfill$\square$\vskip 2ex}
\fi
\expandafter\ifx\csname psmallmatrix\endcsname\relax
\newenvironment{psmallmatrix}
  {\left(\begin{smallmatrix}}
\fi
\expandafter\ifx\csname remark\endcsname\relax
\newenvironment{remark}{{\bf Remark.}}{\hfill\vskip 2ex}
\fi
\newcommand{\spn}{\mathrm{span}}
\makeatother
\title[Rigidity theorems for cone structures]{Rigidity theorems for cone structures}
\author{Tymon Frelik, Wojciech Kry\'nski}
\date{Source: arXiv:2608.12907v1 (CC BY 4.0)}
\begin{document}
\maketitle

\noindent\emph{Source and licence.} Rigidity theorems for cone structures. Version arXiv:2608.12907v1 (CC BY 4.0), distributed under the Creative Commons Attribution 4.0 International licence, as recorded in the arXiv licence metadata for this exact version and shown on the abstract page https://arxiv.org/abs/2608.12907v1. Full rights notice: licenses/arxiv-2608.12907-CC-BY-4.0.txt.\par\smallskip

\noindent\textit{Editorial scope.} Counted unit: the Theorem whose source label is \texttt{thm1} in the article (source0.tex lines 240--242, source label \texttt{thm1}), delivered together with the article's own complete original proof of it (source0.tex lines 243--337, one \texttt{proof} environment of 95 lines). No mathematics is written, repaired or re-derived: statement and proof are the article's own text line for line. Required context retained uncounted: prop1 (lines 221--223). Every definition, display and label of those lines is retained unchanged. Omitted from this package and not redistributed: 1--220, 224--239, 338--760. Nothing inside a retained excerpt is omitted or altered: every retained body is delivered exactly as the article's lines are written, so no omission receipt of any kind accompanies this package. Citations: the retained excerpts carry no citation command at all, so no bibliography entry of the article is transcribed and no reference is redistributed. Reference adaptation: none was needed and none was made; every reference inside the delivered unit resolves inside it, so no unresolved reference or citation is produced. Printed slips kept literal: no printed word, symbol, hypothesis or number of the counted statement or of its proof was altered. Layout and class: the article's own class is \texttt{amsart} with packages including geometry, color, hyperref; the wrapper is the lane's pinned \texttt{amsart} editorial wrapper at 10pt on A4, which restates the theorem-like environments the retained text needs and adds the article's own macro definitions that the retained text needs (\texttt{\textbackslash CC}, \texttt{\textbackslash N}, \texttt{\textbackslash PP}, \texttt{\textbackslash R}, \texttt{\textbackslash Set}, \texttt{\textbackslash VV}, \texttt{\textbackslash XX}, \texttt{\textbackslash dd}, \texttt{\textbackslash spn}), with extra wrapper packages (bm, bbm, dsfont, xspace, cancel, mathtools). The delivered numbering is the unit's own. Assets: no retained span contains a figure environment, a graphics-inclusion command, a PDF-page inclusion, a table or a TikZ picture; no image file is copied into this package, the package declares no asset dependency and no image file is redistributed with it. Licence: version arXiv:2608.12907v1 of this article is distributed under the Creative Commons Attribution 4.0 International licence, as recorded in the arXiv licence metadata for this exact version and shown on the abstract page https://arxiv.org/abs/2608.12907v1. A full rights notice is delivered at licenses/arxiv-2608.12907-CC-BY-4.0.txt. Characters: every retained excerpt is pure ASCII, so the delivered engine, XeLaTeX with its default Latin Modern text font, reports no missing character.
\begin{proposition}\label{prop1}
A non-degenerate cone structure $\CC$ on a manifold $M$ of dimension $n>2$, such that $\CC_x$ is a curve in $\PP(T_xM)$ for any $x\in M$, is of equation type if and only if the associated pair $(\XX,\VV)$ on $N^\CC$ satisfies $[\VV^i,\VV^i]=\VV^{i+1}$ for $i=0,\ldots n-2$.
\end{proposition}

\begin{theorem}\label{thm1}
There exist open dense subsets $\Set_1^{n-1} \subset \mathfrak{V}^{n-1}_1$, for $n>3$, such that any isotrivial cone structure modeled on $\CC_0 \in \Set_1^{n-1}$ arising from a scalar ordinary differential equation is flat and the corresponding equation is linear.
\end{theorem}

\begin{proof}
Let $M$ be a manifold of dimension $n$. Fix an $n$-dimensional vector space $V$ and a basis $(v_0,\ldots, v_{n-1})$ in $V$. We shall identify $\PP^{n-1}$ with $\PP(V)$. Then, there exist functions $\phi^i$, $i=0,\ldots,n-1$, such that an open subset of $\CC_0\subset\PP^{n-1}$ can be parameterized in the following way
\[
\lambda\mapsto \spn\left\{\phi^0(\lambda)v_0+\cdots+\phi^{n-1}(\lambda)v_{n-1}\right\}\in \PP(V),
\]
Let $\CC$ be an isotrivial cone structure modeled on $\CC_0$. It follows from the isotriviality of $\CC$ that any point in $M$ has a neighbourhood with a local frame $(X_0,\ldots,X_{n-1})$ such that, for any $x$ in this neighbourhood, there is a linear map $f_x\colon T_xM\to V$ sending $X_i$ to $v_i$ and transforming  $\hat \CC_x$ to $\hat\CC_0$. By abuse of notation, let us still denote the pullback $f^*\phi$ as $\phi$ for brevity. We get that $\CC_x$ can be locally parameterized as
\[
(x,\lambda)\mapsto \R(\phi^0(\lambda)X_0(x)+\cdots+\phi^{n-1}(\lambda)X_{n-1}(x))\in \PP(T_xM).
\]
From now on we shall assume that $\CC_0$ has no symmetry. It follows that the frames $(X_0,\ldots,X_{n-1})$ as above are only given up to a conformal rescaling:
\[
(X_0,\ldots,X_{n-1})\mapsto \kappa (X_0,\ldots,X_{n-1}),
\]
where $\kappa\colon M\to\R$ is a positive function.


Let $c_{ij}^k$ be the structure functions of the frame $(X_0,\ldots,X_{n-1})$, i.e.,
\[
[X_i,X_j]=\sum_{k=0}^{n-1}c_{ij}^kX_k.
\]
It is a matter of computations to show that the structure functions transform as follows
\begin{equation}\label{conformal}
c_{ij}^k\mapsto\kappa c_{ij}^k+\delta_j^kX_i(\kappa)-\delta_i^kX_j(\kappa)
\end{equation}
under conformal transformations of the frame.

Our goal is to prove that $\smash{c_{ij}^k}=0$ modulo the conformal transformations \eqref{conformal}, for all cones $\CC_0$ from an open and dense subset of $\mathfrak{V}^{n-1}_1$ provided that 
\[
[\VV^i,\VV^i]=\VV^{i+1},
\]
as in Proposition \ref{prop1}. Indeed, if the structure functions $\smash{c^k_{ij}}$ of the frame vanish (modulo conformal transformations), then the frame is locally equivalent (modulo a conformal transformation) to the standard frame in the Euclidean space and consequently the structure is flat.

We shall denote
\[
V(x,\lambda)=\phi^0(\lambda)X_0(x)+\cdots+\phi^{n-1}(\lambda)X_{n-1}(x).
\]
Then,
\[
V_i=V^{(i)}=\left(\tfrac{\dd}{\dd\lambda}\right)^iV=(\phi^0)^{(i)}X_0+\cdots+(\phi^{n-1})^{(i)}X_{n-1}
\]
are consecutive derivatives of $V(x,\lambda)$ with respect to $\lambda$.

Notice that variables $(x,\lambda)$ form a local coordinate system on $N^\CC$. We can consider the vector fields $V_i(x,\lambda)$ as vector fields on $N^\CC$ and we obtain
\[
\begin{aligned}
\XX&=\spn\{\partial_\lambda\}\\
\VV^0_{(x,\lambda)}&=\spn\left\{\partial_\lambda,V(x,\lambda)\right\}\\
\VV^1_{(x,\lambda)}&=\spn\left\{\partial_\lambda,V(x,\lambda),V_1(x,\lambda)\right\}\\
&\vdots\\
\VV^{n-1}_{(x,\lambda)}&=\spn\left\{\partial_\lambda,V(x,\lambda),V_1(x,\lambda),\ldots,V_{n-1}(x,\lambda)\right\}.
\end{aligned}
\]
Consequently, the condition for the cone structure to be of equation type can be expressed as
\[
[V_i,V_j]\in\VV^{j+1}
\]
for $i,j=0,\ldots,n-2$ and $i<j$. Note that for $n=3$ this condition is void (see also the remark following Proposition \ref{prop1}). We have
\[
[V_i,V_j]=\sum_{s,t,k=0}^{n-1}(\phi^s)^{(i)}(\phi^t)^{(j)}c_{st}^kX_k.
\]
It follows from the non-degeneracy assumption that there exist matrix-valued functions $\beta=(\beta_i^j)_{i,j=0,\ldots,n-1}$ depending on the parameter $\lambda$ such that
\[
X_i=\sum_{j=0}^{n-1}\beta_{i}^jV_j.
\]
Indeed, $\smash{\beta=(\beta_i^j)_{i,j=0,\ldots,n-1}}$ is the inverse matrix to $\smash{((\phi^{i})^{(j)})_{i,j=0,\ldots,n-1}}$. We emphasize that the matrices depend on $\lambda$ but not on $x$. In particular, the entries $\smash{\beta_i^j}$, similarly to the functions $\phi^i$, depend solely on the fixed curve $\CC_0\subset \PP(V)$. We get that
\[
[V_i,V_j]=\sum_{s,t,k,l=0}^{n-1}(\phi^s)^{(i)}(\phi^t)^{(j)}\beta_k^lc_{st}^kV_l,
\]
and the condition for the structure to be of equation type reads
\begin{equation}\label{system}
\sum_{s,t,k=0}^{n-1}(\phi^s)^{(i)}(\phi^t)^{(j)}\beta_k^lc_{st}^k=0, \quad l=\max(i,j)+1,\ldots,n-1.
\end{equation}
It is a linear system for the structure functions $c_{st}^k$. Notice that under the conformal transformation \eqref{conformal} the left hand side of \eqref{system} does not change. In fact, we can consider a general transformation
\begin{equation}\label{conformal_general}
\tilde c_{ij}^k=\kappa c_{ij}^k+\delta_j^ka_i-\delta_i^ka_j
\end{equation}
for arbitrary functions $a_i$ and prove that $\tilde c_{ij}^k$ solve \eqref{system} if and only if $c_{ij}^k$ solve \eqref{system}.
Indeed, it is sufficient to substitute $c_{st}^k:=\delta_s^ka_t$ into \eqref{system}. For such a choice we get
\[
\sum_{s,t,k=0}^{n-1}(\phi^s)^{(i)}(\phi^t)^{(j)}\beta_k^lc_{st}^k:=\sum_{t=0}^{n-1}(\phi^t)^{(j)}a_t\sum_{s=0}^{n-1}(\phi^s)^{(i)}\beta^l_s=\sum_{t=0}^{n-1}(\phi^t)^{(j)}a_t\delta^l_i,
\]
which is zero for $l\neq i$, and in particular for $l>i$ as in \eqref{system} (we use above that, by definition, $\beta^l_s$ are coefficients of the inverse matrix to $(\phi^s)^{(i)}$). In the case of a conformal transformation we have $a_i=X_i(\kappa)$. Conversely, if the structure functions of a frame satisfy $c_{ij}^k=\delta_j^ka_i-\delta_i^ka_j$ then we can always locally find $\kappa$ such that $a_i=X_i(\kappa)$. This follows directly from the Jacobi identity which implies the compatibility conditions for differential equations $a_i=X_i(\kappa)$, $i=0,\ldots,n-1$.

The coefficients $c_{ij}^k$ are functions on $M$, i.e. depend on $x$ only, in contrast to the coefficients of system \eqref{system} which are functions of $\lambda$. Hence, the system \eqref{system} depends solely on the curve $\CC_0$ and, furthermore, each equation in \eqref{system} can be differentiated with respect to $\lambda$ arbitrarily many times, producing new equations that are also satisfied by $c_{ij}^k$. The resulting system can be evaluated at a fixed value $\lambda=\lambda_0$ giving a linear system for $c_{ij}^k$ with constant coefficients. The coefficients of this system depend on a high jet of the curve $\CC_0$ at $\lambda_0$ (which can, in fact, be made arbitrarily high by successive differentiation). It is expected that the system is overdetermined for generic curves. In fact, we shall show that for the generic $\CC_0$ the system admits only a trivial solution (modulo \eqref{conformal_general}).

The existence of a nontrivial solution can be expressed in terms of the rank of the system, which is given by a polynomial condition on the coefficients (given by the vanishing of certain determinants). Consequently, this defines a closed subset in the space of curves $\mathfrak{V}^{n-1}_1$. To complete the proof, it remains to show that this subset is of positive codimension. Since it is defined by polynomial conditions it suffices to show that there exists at least one curve satisfying the statement of the theorem.

Chose $n$ numbers, $m_0,m_1,\ldots,m_{n-1}\in\N$, all greater than $n$ and such that $m_i<m_j$ for $i<j$ and all $m_s+m_t-m_k$ as well as $n+m_t-m_0$ are different numbers for $k\neq t,s$. Let $\phi^0(\lambda)=\lambda^n+\lambda^{m_0}$ and $\phi^i(\lambda)=\lambda^{m_{i}}$ for $i>0$. Consider equation \eqref{system} with $l=n-1$. Since functions $\phi^i$ are polynomial in $\lambda$, the left hand side, multiplied by the determinant of matrix $((\phi^i)^{(j)})_{i,j=0,\ldots,n-1}$, which is in the denominator of all $\beta^l_k$, is a polynomial in $\lambda$ as well. Firstly, we shall compute the order of the coefficient next to $c^k_{st}$ as a polynomials in $\lambda$. Let $\mu=m_0+\cdots+m_{n-1}$. Then, after clearing the denominator, $\beta^{n-1}_k$ is a polynomial of order $\mu-m_k-\frac{1}{2}(n-2)(n-1)$, which is an order of the determinant of the matrix $((\phi^i)^{(j)})_{j=0,\ldots,n-2; i=0,\ldots,n-1; i\neq k}$. Consequently, the coefficient next to $c^k_{st}$ is of order
\[
m^k_{st}=\mu-m_k-\frac{1}{2}(n-2)(n-1)+m_s+m_t-i-j,
\]
where $i$ and $j$ are fixed for a given equation in the system \eqref{system}. Note that, due to assumptions on $m_i$'s, coefficients of $c^{k_1}_{s_1t_1}$ and $c^{k_2}_{s_2t_2}$ are of the same order only if $k_1=s_1$, $k_2=s_2$ and $t_1=t_2$ (or with the roles of the lower indices swapped). This automatically gives that \eqref{system} implies that $c^k_{st}$ vanish provided that $k\neq s$ and $k\neq t$, since the corresponding coefficients are independent polynomials in $\lambda$. The remaining coefficients are of the form $c^s_{st}$. For any fixed $t$ we get a system of equations for $c^s_{st}$ where $s=0,\ldots,n-1$, $s\neq t$, obtained by fixing a different value of $i$ in \eqref{system}. We have $n-3$ such equations for each $t$, which are independent. As a solution we get that $c^s_{st}$ are fixed up to 2-parameter freedom (for each given $t$).

One additional equation can be obtained by considering the term $\lambda^n$ in $\phi^0$. It contributes with a monomial of a minimal possible order $\mu+n-m_1-m_k-\frac{1}{2}(n-2)(n-1)$ in the formula for $\beta^{n-1}_k$, for $k>0$. Further on, this term gives a term of order $\mu+n-m_2-\frac{1}{2}(n-2)(n-1)+m_t-i-j$ in $\lambda$ standing next to an expression involving $c^s_{st}$ where $s>0$. This coefficient has to vanish provided that \eqref{system} holds. Ultimately we get, together with the previously obtained $n-3$ equations, $n-2$ independent equations for $c^s_{st}$ (for each $t$), which fix the coefficients uniquely up to conformal transformations \eqref{conformal_general}. This completes the proof.
\end{proof}

\end{document}
